\documentclass[11pt]{article}
\usepackage[T1]{fontenc}
\usepackage{lmodern}
\usepackage[a4paper,margin=25mm]{geometry}
\usepackage{amsmath,amssymb,amsthm}
\usepackage[hypertexnames=false,hidelinks]{hyperref}
\newcommand{\Z}{\mathbb{Z}}
\newcommand{\V}{\mathsf V}
\newcommand{\E}{\mathsf E}
\newcommand{\loc}[1]{\par\smallskip\noindent{\small\textit{Source: #1}}\par\smallskip}
\setlength{\emergencystretch}{2em}
\begin{document}
\begin{center}
{\large Probability on Trees and Networks\\[2pt]Exercise 6.78 and its Author Comment}\\[4pt]
Russell Lyons\quad Yuval Peres\\[2pt]
Version of 19 August 2026---corrected and updated paperback edition
\end{center}

% BEGIN CONTEXT 6.78
\noindent\textbf{Source and scope (editorial).}
This private reader reproduces selected English author text from the fixed 716-page paperback PDF (21,010,205 bytes). Exercise 6.78 is on physical p.\ 251 (printed p.\ 234), in the collected in-text exercises of Chapter 6. Its complete Comment is on physical p.\ 637 (printed p.\ 620). The group that the exercise names is defined in the body on physical p.\ 106 (printed p.\ 89), and its exponential growth is shown on physical p.\ 107 (printed p.\ 90). The upper bound that the exercise says matches its own is Corollary 6.32(iii), printed on physical p.\ 226 (printed p.\ 209); the expansion input behind that corollary is Theorem 6.29 on physical p.\ 224 (printed p.\ 207), followed by the author's remark about the lamplighter group on physical p.\ 225 (printed p.\ 208). The external work that the Comment cites, Pittet and Saloff-Coste (1999), is in the bibliography on physical p.\ 696 (printed p.\ 679). All of this is reprinted below as uncounted context.

The separate errata sheet dated 19 August 2026 contains exactly one Chapter 6 entry, the p.\ 616 correction to Exercise 6.34, and no entry for Exercise 6.78 or for its Comment. The one lamplighter-related entry in that sheet, on p.\ 650, replaces two bibliographic items dated 2015 that are printed on pp.\ 509 and 510; it does not touch the 1999 item cited here, and it does not touch either page of this exercise. Accordingly no revision is applied anywhere in this reader, and no printed text here is altered. No source mathematics is repaired or independently certified.

\medskip\noindent\textbf{Literal source features kept as printed (no synchronization, no repair).}
\begin{itemize}
\item The exercise prints the group as a blackboard-bold capital $\Z$ followed by a script-size circled dot set above the line; the same two-glyph setting is used where the group is defined on physical p.\ 106. It is transcribed here as $\Z^{\odot}$. The letter is taken from a direct full-resolution viewing of physical p.\ 251 recorded in the root session's source inspection of this exercise; it is not derived from the PDF text layer, where the letter is an unmapped byte of the math-italic font (that font carries no Unicode map) and the circled dot comes out as the ordinary circled-dot character at script size.
\item In the PDF text layer of physical p.\ 251 two control characters (bytes 0x0A and 0x0B of the math extension font) sit at the two positions of the fence around the pair $(\Psi_j,X_j)$. They are extraction artefacts of that fence, not printed characters, and nothing is transcribed from them.
\item The exercise prints the lower bound for $p_{2n}(o,o)$ and the Comment prints it for $p_n(o,o)$. On both pages the subscript characters are set at script size on the baseline, not raised above it, so both are ordinary subscripts and neither is an exponent. Neither is synchronized with the other and no index is changed.
\item The exponent $n^{1/3}$ on physical p.\ 251 is raised above the line with the $1$, the solidus and the $3$ all at script size, in contrast with the subscripts above; it is transcribed as printed.
\item The two $j$'s of the pair are set at script size; the printed line breaks inside the exercise are reflowed here and are not content.
\item The constants $c$, $C$ and $c_1$ of the Comment, the ceiling $\lceil n/r^{2}\rceil$ and the interval $[-2r,2r]$ are transcribed as printed. No constant is identified with another and no small-$n$ adjustment is inserted.
\end{itemize}

\medskip\noindent\textbf{Use of the exercise in the body of the work.}
No body theorem, proposition, lemma or example of the fixed version cites Exercise 6.78 by number; the only places in the 716-page PDF where the number occurs are the statement on physical p.\ 251 and the Comment entry on physical p.\ 637. The exercise points outward itself, to Corollary 6.32(iii); part (iii) is what makes Theorem 6.29, through exponential growth, the expansion input for the upper bound it says matches. The remark on physical p.\ 225 records that both sides of Theorem 6.29 are of order $r$ for boxes of radius $r$ in the lamplighter group, so that bound is sharp for this group up to constant factors.

\subsection*{The lamplighter group}
\loc{Physical p.\ 106 (printed p.\ 89).}
The following infinitely presented group will serve as a useful example in several chapters. This is the group $\Z^{\odot}$, also known as the lamplighter group over $\Z$. It is defined as a (restricted) wreath product, which is a special kind of semidirect product. First, the group $\sum_{x\in\Z}\Z_2$, the direct sum of copies of $\Z_2$ indexed by $\Z$, is the group of maps $\psi:\Z\to\Z_2$ with $\psi^{-1}(\{1\})$ finite and with componentwise addition mod 2, which we denote $\oplus$; that is,
\[
(\psi\oplus\psi')(j):=\psi(j)+\psi'(j)\quad(\text{mod }2).
\]
Let $S$ be the left shift, $S(\psi)(j):=\psi(j+1)$. Now define
\[
\Z^{\odot}:=\Bigl(\sum_{x\in\Z}\Z_2\Bigr)\rtimes\Z,
\]
which is the set $\bigl(\sum_{x\in\Z}\Z_2\bigr)\times\Z$ with the following group operation: for $\psi,\psi'\in\sum_{x\in\Z}\Z_2$ and $m,m'\in\Z$, we multiply by the rule
\[
(\psi,m)(\psi',m'):=(\psi\oplus S^{-m}\psi',\,m+m').
\]
We call an element $\psi\in\sum_{m\in\Z}\Z_2$ a configuration and call $\psi(k)$ the bit at $k$. We identify $\Z_2$ with $\{0,1\}$. The second component of an element $x=(\psi,m)\in\Z^{\odot}$ is called the position of the marker in the state $x$. (Another notation for $\Z^{\odot}$ is the restricted wreath product $\Z_2\wr\Z$.) One nice set of generators of $\Z^{\odot}$ is $\{(0,1),(0,-1),(1_{\{0\}},0)\}$. The reason for the name of this group is that we may think of a streetlamp at each integer with the configuration $\psi$ representing which lights are on, namely, those where $\psi=1$. We also may imagine a lamplighter at the position of the marker. The first two generators of $\Z^{\odot}$ correspond (for right multiplication) to the lamplighter taking a step either to the right or to the left (leaving the lights unchanged); the third generator corresponds to flipping the light at the position of the lamplighter. See Figure 3.4.

\subsection*{Exponential growth of the lamplighter group}
\loc{Physical pp.\ 106--107 (printed pp.\ 89--90).}
The lamplighter group $\Z^{\odot}$ has exponential growth: Consider the subset $T_n$ of group elements at distance $n$ from the identity that can be arrived at from the identity by using only the generators $(0,1)$ and $(1_{\{0\}},0)$, in other words, by never allowing the lamplighter to move leftward. Since $T_n$ is a disjoint union of $(0,1)T_{n-1}$ and $(1_{\{0\}},0)(0,1)T_{n-2}$, we have
\[
|T_n|=|T_{n-1}|+|T_{n-2}|.
\]
Thus, $\{|T_n|\,;\,n\ge0\}$ is the sequence of Fibonacci numbers. Therefore the exponential growth rate of $|T_n|$ equals the golden mean, $(1+\sqrt5)/2$. In fact, it is easy to check that this is the growth rate of balls in $\Z^{\odot}$.

\subsection*{Theorem 6.29 and the expansion input}
\loc{Physical pp.\ 224--225 (printed pp.\ 207--208).}
Define the inner vertex boundary of a set $K$ as
\[
\partial^{in}_{\V}K:=\{x\in K\,;\,\exists y\notin K\ \ y\sim x\}.
\]
\textbf{Theorem 6.29.} (Expansion of Cayley Graphs) \textit{Let $G$ be a Cayley graph. Let $R(m)$ be the smallest radius of a ball in $G$ that contains at least $m$ vertices. Then, for all finite $K\subset\V$, we have}
\[
\frac{|\partial^{in}_{\V}K|}{|K|}\ \ge\ \frac{1}{2R(2|K|)}\,.
\]
The same proof yields the same lower bound for the outer vertex boundary, $\partial_{\V}K$. An extension to transitive graphs is given in Lemma 10.46.

Note that both sides of Theorem 6.29 are of order $1/r$ for balls of radius $r$ in $\Z^d$ ($d\ge1$) and also of order $r$ for ``boxes'' (6.8) of ``radius'' $r$ in the lamplighter group $\Z^{\odot}$, so, up to constant factors, this bound is sharp for such groups. Of course, it is not sharp for nonamenable groups.

\subsection*{Corollary 6.32: the bound the exercise names}
\loc{Physical pp.\ 225--226 (printed pp.\ 208--209).}
For $S,A\subset\V$, let $Q(S,A):=\sum_{s\in S,a\in A}Q(s,a)$ and $|\partial_{\E}S|_Q:=Q(S,S^c)$. The edge expansion of a finite set $S\subset\V$ is $\Phi_S:=|\partial_{\E}S|_Q/\pi(S)$, and the expansion profile of the chain is defined for $u\in(\pi_{\min},\infty)$ by
\[
\Phi(u):=\inf\{\Phi_S\,;\,0<\pi(S)\le u\}.
\]

\textbf{Corollary 6.32.} \textit{Suppose that the Markov chain $(\V,P)$ is lazy or reversible.}
\begin{enumerate}
\item[(i)] If $\Phi(u)\ge\varphi_0$ for some $\varphi_0>0$ and all $u>0$, then
\[
\frac{p_n(x,y)}{\pi(y)}\ \le\ \frac{4}{\pi_{\min}}\exp\Bigl(-\frac{\varphi_0^{2}(n-1)}{16}\Bigr)
\qquad\text{for all }x,y\in\V\text{ and }n\ge1.
\]
\item[(ii)] Let $d>0$. If $\Phi(u)\ge cu^{-1/d}$ for some $c>0$ and all $u>0$, then
\[
\frac{p_n(x,y)}{\pi(y)}\ \le\ C'n^{-d/2}
\qquad\text{for all }x,y\in\V\text{ and }n\ge1,
\]
where $C'=C'(d,c)$. In particular, by Theorem 6.29, this applies if the chain is simple random walk on a Cayley graph where the balls satisfy $|B(o,r)|\ge c_0r^d$ for some $c_0>0$, so $R(k)=O(k^{1/d})$.
\item[(iii)] If $\Phi(u)\ge c/\log(bu)$ for some $b,c>0$ and all $u>0$, then
\[
\frac{p_n(x,y)}{\pi(y)}\ \le\ C_1\exp(-C_2n^{1/3})
\qquad\text{for all }x,y\in\V\text{ and }n\ge1,
\]
where $C_1$, $C_2$ depend on $b$, $c$. In particular, this applies if the chain is simple random walk on a Cayley graph of exponential growth because of Theorem 6.29.
\end{enumerate}
% END CONTEXT 6.78

\clearpage
% BEGIN PRIMARY STATEMENT 6.78
\section*{Exercise 6.78}
\loc{Physical p.\ 251 (printed p.\ 234).}
\begin{sloppypar}
\noindent\textbf{Exercise 6.78.} Show that simple random walk $(\Psi_j,X_j)$ on the Cayley graph of $\Z^{\odot}$ (with the standard generators) satisfies $p_{2n}(o,o)\ge\exp(-Cn^{1/3})$ for some $C$ and all $n$. (The matching upper bound was proved in Corollary 6.32(iii).)
\end{sloppypar}
% END PRIMARY STATEMENT 6.78

% BEGIN PRIMARY PROOF 6.78
\section*{Complete author Comment on Exercise 6.78}
\loc{Physical p.\ 637 (printed p.\ 620).}
\textbf{6.78.} For every $v\in[-r,r]$, we have
\[
\mathsf P_v\Bigl[\max_{j\le r^{2}}|X_j|\le 2r\ \text{ and }\ |X_{r^{2}}|\le r\Bigr]\ \ge\ e^{-c}
\]
for some constant $c$. Repeatedly using the Markov property at times that are multiples of $r^{2}$, we deduce that
\[
\mathsf P_v\Bigl[\max_{j\le 2n}|X_j|\le 2r\Bigr]\ \ge\ e^{-c\lceil n/r^{2}\rceil}.
\]
Adding the requirement that each lamp in $[-2r,2r]$ is turned off at the last visit to it, and that the lamplighter returns to the origin, we infer that
\[
p_n(o,o)\ \ge\ \exp(-cn/r^{2}-c_1r).
\]
Optimizing this lower bound suggests taking $r=n^{1/3}$. See Pittet and Saloff-Coste (1999).
% END PRIMARY PROOF 6.78

\medskip\noindent\textbf{Cached bibliographic attribution (not an additional source proof).}
The Comment on Exercise 6.78 closes with ``See Pittet and Saloff-Coste (1999).'' The bibliography of the fixed version, on physical p.\ 696 (printed p.\ 679), carries the entry:

\smallskip
\noindent Pittet, C. and Saloff-Coste, L. (1999) Amenable groups, isoperimetric profiles and random walks. In \textit{Geometric Group Theory Down Under} (Canberra, 1996), pages 293--316. de Gruyter, Berlin. MR: 1714851.
\smallskip

\noindent That work is not reproduced, obtained or consulted here; the entry is cached from the book's own bibliography, and the attribution is the author's.
\end{document}
